\pdfoutput=1
\pdfmapfile{+symbols.map}
\pdfmapfile{+cm.map}
\pdfmapfile{+cmextra.map}
\documentclass[10pt]{article}
\usepackage[margin=1in]{geometry}
\usepackage{amsmath,amssymb,amsthm}
\usepackage[hidelinks]{hyperref}
\usepackage{microtype}
\newtheorem{theorem}{Theorem}
\newtheorem{lemma}{Lemma}
\title{A counterexample to the core partition component formula}
\author{TLMC conjecture 00000003829}
\date{11 October 2026}
\renewcommand{\labelitemi}{$\bullet$}
\begin{document}
\emergencystretch=3em
\maketitle
\begin{abstract}
The coprime component formula in conjecture 00000003829 fails at
$(s,t)=(2,3)$. There are exactly two simultaneous cores, the empty
partition and the one-box partition. A single-box border joins them,
so the graph has one connected component, whereas the asserted rational
Catalan value is two. The complete classification, the border, and the
component-count contradiction are verified in a self-contained Lean 4
project using ordinary kernel-checked proofs.
\end{abstract}

\section{Statement and border convention}
The conjecture defines $B_{s,t}$ as the graph of $(s,t)$-core partitions
generated by adding and removing borders, and asserts for coprime $s,t$
that
\[
\#\pi_0(B_{s,t})=\frac{1}{s+t}\binom{s+t}{s}.
\]
A partition is a finite weakly decreasing sequence of positive integers;
the empty partition is included. Its diagram has $\lambda_i$ cells in
row $i$. The hook length is the number of cells strictly right of a
cell, plus those strictly below it, plus one. A $t$-core has no hook
length divisible by $t$.

We use the usual border-strip (rim-hook or ribbon) convention: a
nonempty skew diagram connecting two partitions, connected through
shared cell edges and containing no $2\times2$ square. In particular,
a single cell is a border strip. These standard definitions appear in
Jouhet and Wahiche, page~2~\cite{jw}. The same counterexample applies to
the single-box Young graph, and to any border relation that retains the
single-box edge used below.

The filed English and Chinese statements impose no restriction on border
length. If a different graph is intended, for example one allowing only
borders of lengths $s$ or $t$, its edges need to be specified separately;
the present counterexample does not assert the same conclusion for that
modified definition. The Lean project explicitly records its graph
rather than leaving the adjacency convention implicit.

\section{Classification without a size bound}
\begin{lemma}
The simultaneous $(2,3)$-core partitions are exactly
$\varnothing$ and $(1)$.
\end{lemma}
\begin{proof}
Both displayed diagrams have no hook divisible by $2$ or $3$.
For the converse, induct on the number of rows. Removing the first row
does not change any hook in the lower rows, so the remaining partition
is itself a simultaneous core. By induction it is empty or $(1)$.

If it is empty, the original partition is one row of length $x>0$.
If $x\ge2$, its second cell from the right has hook length $2$,
a contradiction. Thus $x=1$.

If the remaining partition is $(1)$, the original partition has row
lengths $(x,1)$. When $x=1$, the top-left cell has hook length $2$.
When $x=2$, it has hook length $3$. When $x\ge3$, the second cell
from the right in the first row has no cell below it and has hook length
$2$. Every possibility contradicts a core condition. Thus this case
cannot occur, completing the induction.
\end{proof}

\newpage
\begin{theorem}
Under the stated border convention, conjecture 00000003829 is false.
\end{theorem}
\begin{proof}
Take $s=2$ and $t=3$. The graph has precisely the two vertices from the
lemma. Adding the cell $(0,0)$ to the empty diagram produces $(1)$.
The skew diagram consists of this one cell, so it is nonempty,
edge-connected, and contains no $2\times2$ square. Hence the two vertices
are adjacent and their component quotient has cardinality one.
On the other hand,
\[
\gcd(2,3)=1,\qquad
\frac{1}{2+3}\binom{2+3}{2}=\frac{10}{5}=2.
\]
The component count is therefore different from the claimed value.
The coprime clause already fails, irrespective of the conjecture's
additional noncoprime clause.
\end{proof}

\section{Formal verification and reproducibility}
The file \texttt{Basic.lean} imports only \texttt{Std} and is pinned to
Lean \texttt{4.31.0}. It does not impose an enumeration cutoff on
partitions. Its main definitions and theorems are:
\begin{itemize}
\item \texttt{hook}, \texttt{IsPartition}, and \texttt{IsCore} encode
ordinary diagrams and every hook in every row.
\item \texttt{classify} proves the lemma for arbitrary finite lists of
positive row lengths satisfying both core conditions. This is slightly
stronger than needed, since the proof does not use row monotonicity.
\item \texttt{cells}, \texttt{skew}, and \texttt{Border} encode actual
Young diagrams and border strips; \texttt{singleton\_border} verifies
the required edge.
\item \texttt{Vertex} contains all simultaneous cores. \texttt{Reach}
is undirected path equivalence and \texttt{Component} is its quotient.
\texttt{HasCard} asserts the existence of a bijection with
$\mathrm{Fin}(n)$.
\item The quotient has one element (\texttt{component\_count\_one})
and cannot have two (\texttt{component\_count\_not\_two}).
\item \texttt{connected\_for\_any\_border\_relation} proves that any
edge relation on these vertices containing the singleton edge is
connected. Thus additional allowed borders cannot change the result.
\item \texttt{counterexample} combines coprimality, cardinality one,
and the failure of the conjectured value. Pascal recursion calculates
$\binom52=10$ and $10/5=2$ exactly; no rounding affects the witness.
\end{itemize}

With the pinned toolchain on the executable search path, run
\texttt{./reproduce.sh}. It builds the project, re-elaborates the proof,
and replays its declarations through \texttt{leanchecker}.
The axiom audit for the final theorem contains only Lean's standard
\texttt{propext} and \texttt{Quot.sound}. No unproved assumptions,
external computation oracle, or kernel bypass is used. The supplied
\texttt{build.log} records the successful run.

\begin{thebibliography}{9}
\bibitem{problem} The Last Math Competition.
\emph{Conjecture 00000003829}, retrieved 11 October 2026.
\href{https://github.com/The-Last-Math-Competition/The-Last-Math-Competition/blob/9b795e7a94a6076e49a65a6489e1caf9153abd23/conjectures/00000003829.md}{Official statement at the checked upstream commit}.
\bibitem{jw} Fr\'ed\'eric Jouhet and David Wahiche.
\emph{Congruences for hook lengths of partitions}, page 2.
\url{https://math.univ-lyon1.fr/~jouhet/congruences_pams_vf.pdf}.
\end{thebibliography}
\end{document}
